% Ad Atticum 9.16 (ad-atticum-09-16) --- English translation
% Translated by Claude (Anthropic) from the Latin (Perseus).
% Style guide: STYLE.md.

\ciceroLetterOpener{CICERO ATTICO salutem}

\ciceroSection{1} Though I had nothing to write to you, still, that I might not let a day slip by, I send this letter. They report that Caesar will lodge on the sixth day before the Kalends at Sinuessa. From him a letter was put into my hand on the seventh day before the Kalends, in which he is now looking for my resources, not, as in the earlier letter, for my help. When I had praised by letter that famous clemency of his at Corfinium, he wrote back in these terms:

\ciceroSection{2} ``Caesar \textit{imperator} to Cicero \textit{imperator}, greeting. You divine rightly about me --- for you know me well --- that nothing lies further from me than cruelty. And as I myself take great pleasure in the thing itself, so I also rejoice in triumph that my conduct meets with your approval. Nor does it move me that those whom I let go are said to have made off in order to bring war upon me again. I want nothing more than that I should be like myself, and they like themselves.''

\ciceroSection{3} I should be glad if you would be at hand to me near the City, so that I may, as I have been used to, draw on your counsel and your good offices in everything. Of your Dolabella, you may be sure, nothing is more pleasing to me. Indeed I shall owe him this thanks for it; for he cannot do otherwise --- such is his humanity, such his good sense, such his goodwill toward me.
